%% kapa: exam=1 round=20 year=2009 period=1 form=A subject=economics subject_ko=경제원론 seq=8 number=48
%% source: Q-Net 감정평가사 기출문제 (cst00309) artlSeq=5024178 file=extracted/1008/5024178/2009년 제20회 감정평가사 1차시험 기출문제/1교시.hwp
%% conversion: kordoc 4.24.51 -> tools/gen1.py (auto); review=proofread-sonnet; flags=tikz
\begin{question}{48}
자동차 등록세는 자동차 구입가격의 7\% 이다. 정부는 신차를 구입하는 사람들에게 등록세를 70\% 감면해주는 정책을 발표했다. 세금 감면을 실시한 후 예상되는 자동차 수요곡선은? (단, 그림에서 굵은 실선은 정책 시행 이전의 수요곡선이고, 점선들은 정책 시행 이후의 가상 수요곡선들이다.)
\begin{center}
\begin{tikzpicture}[font=\small, x=1cm, y=1cm]
  \draw (0,0) -- (0,5.3);
  \draw (0,0) -- (5.5,0);
  \node[above] at (0.2,5.35) {자동차 가격};
  \node[right] at (5.55,0) {자동차 수량};
  \node[below left] at (0,0) {0};
  \draw[line width=1.2pt] (0,2.9) -- (2.87,0);
  \draw[dashed, thin] (0,4.8) -- (4.72,0);
  \draw[dashed, thin] (0,4.8) -- (2.87,0);
  \draw[dashed, thin] (0,2.9) -- (4.72,0);
  \draw[dashed, thin] (0,1.5) -- (2.87,0);
  \draw[dashed, thin] (0,1.5) -- (1.5,0);
  \node at (2.1,2.8) {a};
  \node at (1.1,3.1) {b};
  \node at (2.7,1.25) {c};
  \node at (1.4,0.8) {d};
  \node at (0.8,0.65) {e};
\end{tikzpicture}
\end{center}
\begin{choices}
  \item a
  \item b
  \item c
  \item d
  \item e
\end{choices}
\end{question}
\answer{2}
